\chapter{The polity}
\label{ch:polity}

\section{The problem is irrevocability, not policy}

Chapter~\ref{ch:ledgerch} ended on a finding, and this chapter is its
consequence.

Almost every measure this book proposes has already been tried in Cuba. Private
restaurants: tried, 1993, worked, restricted 2017. Farmers' markets: tried,
1980, worked, abolished 1986, tried again 1994. Dollars: legalised 1993,
withdrawn 2004. Land in usufruct: 2008 and 2012. Private firms with legal
personality: 2021, then price-capped and blamed for inflation by 2024. The
Cuban state has not been short of good ideas. It has been unable to keep any of
them for longer than the emergency that produced them.

The consequence is not merely that reforms get reversed. It is that they
underperform \emph{while in force}, because everyone can see they are
revocable. A Cuban who is deciding whether to invest three years' savings in
equipment is making a bet on the durability of a decree, and the historical
base rate on that bet is poor. So capital is not committed, businesses stay
small and liquid, nobody trains anyone, nobody builds anything that cannot be
moved --- and the resulting mediocrity is then produced as evidence that the
private sector does not work, which justifies the next restriction.

This is a solved problem in institutional economics and the solution has a
name \citep{northweingast1989}. What a state needs, when it wants people to invest behind its promises, is
a mechanism by which it can \emph{disable its own future discretion} in a way
outsiders can verify. Not a better promise --- a promise it cannot cheaply
break. Everything in this chapter is an instrument of that kind, and everything
in the rest of Part III depends on at least some of them existing. A chapter of
excellent energy policy in a country where policy is revocable at will is a
chapter about what could have been.

\begin{design}{Commitment before content}
\label{d:commit}
\begin{rules}
\rl{1}{Ordering} No liberalising measure in Part~III is adopted without at
least one commitment device attached to it: a constitutional entrenchment, a
justiciable right, a treaty obligation, an independent regulator with secured
tenure and budget, or a published rule with an automatic trigger. A measure
adopted by decree alone is understood by everyone, correctly, as temporary.

\rl{2}{Rights not lists} Economic activity is permitted by right, subject to
published prohibitions, and never by enumerated list of allowed categories. The
list is what made the Cuban private sector a collection of trades with no
software developers in it, and it is also the instrument by which the sector is
shrunk without any law being repealed.

\rl{3}{Asymmetry} Instruments are designed so that expanding a freedom is
administratively cheap and withdrawing it is expensive --- requiring
legislation, publication, notice periods and compensation --- rather than the
reverse, which is the present arrangement.

\rl{4}{Standing} Every entitlement created in Part~III is enforceable by the
person who holds it, in a court, against the state. An entitlement nobody can
sue over is a policy, and policies are what this chapter is about not relying
on.
\end{rules}
\end{design}

\section{Sequence: law before ballots}

The most consequential choice in this chapter is the order, and the book takes a
position that will not please everybody.

\emph{Rule of law first; competitive national elections later, and on a
published schedule.}

The reasoning is empirical rather than principled. The post-communist
transitions that went well --- Estonia, Poland, Slovenia, the Czech lands ---
built courts, property registries, competition authorities and independent
central banks alongside or ahead of their electoral opening. The ones that went
badly held elections into an institutional vacuum, where the winner inherited
undivided control of unpriced assets and no court capable of restraining
anybody, and the resulting concentration then captured the electoral process
itself. Russia's 1996 election was won by an incumbent financed by the men to
whom he had sold the economy; the sequence, not the ballot, was the problem
\citep{aslund2013}.

Cuba's own history says the same thing from the other side. Chapter~\ref{ch:inheritance}
argued that the decisive feature of January 1959 was that Batista's coup had
already destroyed every institution capable of arbitrating between contenders,
so that the only body with legitimacy was the guerrilla movement and there was
nothing to slow its consolidation. A Cuban transition that produces one
legitimate actor and no others will end where 1959 ended, whatever its
intentions. The institutions have to exist before the contest, or the contest
decides who owns them.

Two things must be said against this position, because it is the position under
which authoritarian governments always promise elections later.

First, ``later'' has to be a date. A sequencing argument without a published
calendar is an indefinite postponement with a rationale, and the Cuban
government has been offering exactly that since 1959. The commitment device
here is that the calendar is entrenched at the beginning, along with the
institutions, and is itself justiciable.

Second, some political opening is a \emph{precondition} for the institutions
rather than a reward for them. A judiciary cannot be made independent in secret;
an anti-corruption body cannot function without a press to publish what it
finds; a property registry is only trusted if its contents can be disputed in
public. So the sequence is not law-then-politics but a specific ordering
\emph{within} the political opening: association, press, information and courts
early, because the institutions need them; competitive national elections on a
fixed later date, because the assets need pricing and the parties need building
first.

\begin{design}{The political calendar}
\label{d:calendar}
\begin{rules}
\rl{1}{Immediate} Legal association, independent press and publishing,
independent trade unions, and the right to receive and impart information
across borders. These are preconditions for every institution below and cost
the state nothing to concede.

\rl{2}{Year one} Judicial reform (Design~\ref{d:courts}), an administrative
court in which a citizen or firm can sue the state, publication of the state's
accounts, and a party law setting the terms on which political organisations
register --- deliberately permissive, with public financing and disclosure.

\rl{3}{Years two and three} Competitive municipal elections, held first and
twice, because municipalities are where parties learn to be parties, where the
electoral administration learns its job, and where a mistake is survivable.

\rl{4}{Year five} Competitive elections to the national assembly, on a date
fixed in the constitutional amendment that begins the process, with an
independent electoral commission constituted in year one and no incumbent
majority on it.

\rl{5}{Non-postponement} The dates are entrenched and their postponement
requires a supermajority plus judicial confirmation of a stated emergency,
with automatic expiry. A calendar a government can move at will is not a
calendar.
\end{rules}
\end{design}

\section{Courts}

Cuba has had no court that ruled against the executive on a political matter
since 1959, and the practice of subordinating the judiciary to political
direction was established in the first weeks of that year, when acquitted
defendants were retried on the leader's public insistence. Sixty-seven years of
that produces a bench with technical competence in ordinary civil and criminal
work and no institutional memory whatever of independence.

The reform therefore cannot be a decree declaring judges independent. It has to
change the four things that actually determine whether a judge can rule against
the government: who appoints them, how long they serve, who pays them, and who
can remove them.

\begin{design}{Judicial independence}
\label{d:courts}
\begin{rules}
\rl{1}{Appointment} Judges are appointed by a judicial council whose membership
is majority-judicial and whose non-judicial members are appointed by qualified
majority of the assembly, not by the executive. Vacancies are advertised;
criteria and shortlists are published.

\rl{2}{Tenure} Appointment to retirement age, removable only for incapacity or
proven misconduct, by a process that is itself judicial and published. The
present arrangement, in which judges serve at the pleasure of the assembly
that elects them, is abolished.

\rl{3}{Budget} The judiciary's budget is a fixed minimum share of state
spending, set in the constitution, disbursed on the judiciary's own authority
and audited publicly. A judiciary whose electricity bill is paid at the
executive's discretion is not independent, and in Cuba this is not a
metaphor.

\rl{4}{Administrative court} A specialised administrative jurisdiction is
created in which any person or firm may sue any organ of the state over a
licence refused, a permit delayed, a good seized, a contract broken or a
regulation misapplied, with the state bearing costs when it loses. This is the
single most important court in Part~III: it is where the commitment devices of
Design~\ref{d:commit} become real, and it is the mechanism by which
Chapter~\ref{ch:integrity}'s removal of discretion is enforced.

\rl{5}{Publication} All judgments are published in full, in a free public
database, within thirty days. A body of published reasoning is what turns a
set of decisions into a law, and it is also the cheapest available check on
the decisions themselves.

\rl{6}{Commercial capacity} Commercial, insolvency and administrative benches
are staffed by judges trained for them --- a specific, funded, decade-long
programme, since Chapter~\ref{ch:socialstate} established that Cuba trained
very few lawyers to think about how a market works. This is a training
constraint, not a legislative one, and it binds.
\end{rules}
\end{design}

\section{What has to be legal}

The following is a floor rather than an ideal, stated as the minimum on which
the rest of Part III depends. Each item is included because something later
fails without it, and the reason is given.

\begin{design}{The minimum conditions}
\label{d:minimum}
\begin{rules}
\rl{1}{Association} The right to form and join organisations --- parties,
unions, professional bodies, cooperatives, churches, neighbourhood groups ---
by registration on published criteria, with refusal appealable to the
administrative court. Required by: every collective actor in Part~III,
including the farmers' associations of Chapter~\ref{ch:land} and the
professional bodies that must accredit Chapter~\ref{ch:life}'s regulators.

\rl{2}{Press} Independent publishing and broadcasting, by registration, with no
prior restraint and with defamation as a civil rather than a criminal matter.
Required by: Chapter~\ref{ch:integrity}, whose entire enforcement model is that
somebody publishes what the auditor finds. Law 88 of 1999 and the criminal
provisions on enemy propaganda are repealed.

\rl{3}{Information} A statutory right of access to state-held information, with
a presumption of disclosure, a short statutory deadline, a narrow and
enumerated exemption list, and appeal to the administrative court. Required by:
open contracting, asset declarations, the published cadastre, the electoral
register, and the environmental data of Chapter~\ref{ch:nature}.

\rl{4}{Movement} The right to leave and to return, entrenched, with the
remaining professional exit restrictions repealed. Required by:
Chapter~\ref{ch:life}'s reform of medical missions, and by the simple fact that
a state which can prevent departure never has to earn its citizens' consent
to stay.

\rl{5}{Communication} No interruption of internet or telephone service on
political grounds, with any interruption requiring judicial authorisation,
publication and a stated duration. Required by: Chapter~\ref{ch:talent}, whose
entire proposition is that a foreign client can rely on a Cuban supplier being
reachable. The national shutdown of 11 July 2021 would, under this rule, be
unlawful --- and that is the point of writing it down.
\end{rules}
\end{design}

\section{The bargain with the incumbents}

Chapter~\ref{ch:premises} established that the party, the armed forces and
GAESA are parties to any transition rather than obstacles to be removed from
one. This section states the bargain, because a bargain that is left implicit is
one that gets renegotiated at every step.

The exchange is: \emph{security, standing and property, in return for the
surrender of discretion.} The incumbents keep their liberty, their pensions,
their homes, a legal route into politics, and --- within limits and through a
transparent process --- an interest in the assets they now run. They give up the
power to decide cases, allocate licences, set exchange rates, direct
prosecutions and control foreign currency.

This is not generosity and it should not be presented as such. It is the price
of a transition that does not go through a collapse, and the alternative
histories are available: the transitions that offered the outgoing elite nothing
were the violent ones, and the ones that offered them everything produced
oligarchies. The design's job is to be specific about where the line falls.

\begin{design}{Amnesty, and where it stops}
\label{d:amnesty}
\begin{rules}
\rl{1}{Covered} Amnesty for the fact of having held office, belonged to the
party, served in the security organs, administered the economy, enforced the
laws then in force, or participated in the political system as it existed.
Nobody is prosecuted for having been part of it.

\rl{2}{Covered} Amnesty for economic offences below a published threshold, and
for offences arising from the ordinary conduct of a state enterprise under the
rules then applying.

\rl{3}{Not covered} Killing, torture, enforced disappearance, sexual violence,
and orders to commit them. These are not amnestied, at any rank, and no statute
of limitations runs on them. A transition that amnesties torture teaches that
torture is a policy option.

\rl{4}{Not covered} Large-scale personal enrichment through the abuse of
office, above the threshold in clause~2, where the asset can be traced. The
remedy here is recovery of the asset rather than imprisonment, which is both
more achievable and more useful.

\rl{5}{Conditional} Amnesty under clauses~1 and~2 is conditional on full
disclosure to the commission in Design~\ref{d:truth} and on the surrender
of undeclared foreign assets. It is revocable for a materially false
declaration, and only for that.

\rl{6}{Symmetric} The amnesty runs in both directions. Everyone convicted for
political expression, association, unauthorised departure, or the events of 11
July 2021 is released and their conviction vacated. This is not a concession
traded against the others; it is a precondition, and it happens first.
\end{rules}
\end{design}

\begin{design}{Truth without prosecution}
\label{d:truth}
\begin{rules}
\rl{1}{Mandate} A commission with a fixed term of three years, a duty to take
testimony from any person, powers to compel documents from the state, and a
duty to publish. It has no power to prosecute, to punish, or to name a person
as criminally liable.

\rl{2}{Scope} The whole period, both directions, and both kinds of harm: the
executions and prison terms, the camps, the purges, the confiscations and the
repudiation mobs; and also the deaths caused by the exile organisations, the
sabotage campaigns, and the Bay of Pigs. A commission that examines one side
produces a document that half the country will not read.

\rl{3}{Records} The state archives on these matters are transferred to the
commission and thereafter to a public archive, with a fixed and short closure
period for genuine personal privacy and nothing else.

\rl{4}{Restitution of standing} The commission may recommend, and the state
must act on, the restoration of professional standing, pensions and academic
records to people removed under \emph{parametraci\'on} between 1971 and 1976
and to their equivalents in later periods. Much of that harm was done by
administrative letter and can be undone the same way, at almost no cost, and it
is the cheapest act of repair available.

\rl{5}{No hierarchy of victims} The commission does not rank suffering, adjudicate
between narratives, or produce a verdict on the revolution. It produces a record.
\end{rules}
\end{design}

\section{Where power sits}

One structural point, easy to omit and expensive to omit.

Cuban centralism is older than the revolution. The republic was run from Havana,
the colony was run from Havana, and the 1976 constitution's \emph{Poder Popular}
gave municipalities elections without giving them money. A municipality that
cannot raise or retain revenue cannot maintain a road, a clinic or a water
pump, and so the entire country's problems arrive at the same desk in the
capital, where they queue.

This matters for Part III specifically because so much of what it proposes is
local: rooftop generation and distribution in Chapter~\ref{ch:energy}, licensing
and food procurement in Chapter~\ref{ch:tourism}, the cadastre and building
control in Chapter~\ref{ch:capital}, coastal setbacks in
Chapter~\ref{ch:nature}, and the property tax in Chapter~\ref{ch:welfare}, which
is the natural local revenue and the instrument that makes local government
answerable to local residents rather than to the ministry.

\begin{design}{Municipal capacity}
\label{d:municipal}
\begin{rules}
\rl{1}{Own revenue} Municipalities retain the property tax and a fixed share of
the accommodation levy and of business licence fees raised in their territory,
set by formula and not by negotiation.

\rl{2}{Equalisation} A transfer formula, published and rule-based, redistributes
toward the eastern provinces on measured indicators of need. Without this,
Chapter~\ref{ch:premises}'s regional constraint becomes a regional divergence:
tourism revenue accrues where the tourists are, and the tourists are not in
Guant\'anamo.

\rl{3}{Competence} Municipalities hold and exercise building control, local
licensing, local roads, water and waste, and primary-care and school premises,
with the service standards set nationally and the delivery local.

\rl{4}{Accounts} Every municipality publishes its budget, its execution and
its contracts, in a common machine-readable format, quarterly. This is also
Chapter~\ref{ch:integrity}'s cheapest and most effective instrument, because
local theft is visible to local people if the numbers are available to them.

\rl{5}{Internal movement} Decree 217 of 1997 and the associated restrictions on
residence in Havana are repealed. Restricting internal migration does not
develop the east; it just makes being from the east a legal disability.
\end{rules}
\end{design}

\section{What this chapter does not settle}

Two things are left open, deliberately, and a reader should know which.

The first is the form of the state --- presidential or parliamentary, unicameral
or bicameral, the electoral formula. These matter, they are extensively studied
elsewhere, and nothing in this book's economics turns on them. They belong to
whoever writes the constitution and should be settled by Cubans in a
constituent process rather than proposed by an economist in a chapter.

The second is who begins. The design in this chapter is executable by the
present government --- every clause of it could be enacted by the National
Assembly as presently constituted, and several of them, notably the amnesty in
both directions and the release of communications from political interruption,
would cost it very little and buy it a great deal. It is also executable by a
successor government. It is not executable by nobody, and
Chapter~\ref{ch:sequence} returns to the question of which coalition carries
which phase.
